% ---------------------------------------------------------------------------
% Results of the IMERG/BMD ingestion screen.
%
% DRAFTING NOTE: as with docs/ablation_v3.tex this is deliberately long, so
% that it can be cut down later without re-deriving anything. Source data:
%   data/processed/ing2022_report/ingestion_matrix.{csv,json}
%   data/processed/ing2022_report/selection_{background,gauges,satellite,
%                                            combined}/config_selection.json
%   data/processed/bmd_imerg_eval_ing2022_*/fold[0-4].npz
% Design: docs/EXPERIMENT_ingestion.md
% ---------------------------------------------------------------------------

\section{How should IMERG and BMD gauges be combined?}
\label{sec:ingestion}

\subsection{Design}

Eleven assimilation configurations were run over 2022-05-01 to 2022-05-10 with
five rotated station folds and 30 ensemble members, giving 380 withheld
station-days of which 178 were wet ($\geq 1$~mm~day$^{-1}$). The prior was held
fixed at V1 (\texttt{runs/prior\_h100\_cpc/best.pt}); nothing here retrains
anything. Every run writes four arms from matched seeds --- background,
gauges-only, satellite-only and simultaneous --- so no separate gauges-only
configuration is required.

The arms vary along two axes. \emph{Scale}: the same IMERG field assimilated at
$0.1^\circ$, $0.4^\circ$ and $0.8^\circ$ footprints. \emph{Strength}: gauge
$\sigma_{\mathrm{obs}}$ at $0.05$, $0.25$ and $0.41$; satellite
$\sigma_{\mathrm{obs}}$ at $0.20$ and $1.00$; their ratio; and the measured
observation error from the variogram analysis. One thinning contrast (stride 1
against the default stride 3) tests whether an earlier null transfers.

Two decisions constrain the design and are recorded because they are not
obvious.

\paragraph{The $R$ multiplier is not a usable strength knob.}
\texttt{perturb\_observations} draws each member's observation perturbation with
standard deviation $\sqrt{R}$, so inflating $R$ inflates the perturbations
rather than merely down-weighting the observation. Strength is therefore varied
through $\sigma_{\mathrm{obs}}$ throughout, which reaches the same effective
weight without injecting noise.

\paragraph{Observation scale requires a new observation file.}
\texttt{observations.imerg.factor} declares which grid the observation file
lies on: the assimilation constructs the expected footprint centres from the
model grid and requires the file's coordinates to match them. Changing the
factor alone is therefore not a scale experiment but an error. The coarse arms
use block-averaged observation files
(\texttt{scripts/44\_coarsen\_imerg\_observations.py}), with three settings
moved in step: the file, the factor, and \texttt{error\_corr\_cells}, which is
denominated in cells \emph{of the observation grid} and must be rescaled from
$3.0$ to $0.75$ and $0.375$ to hold the physical correlation length fixed at
$33$~km. Thinning is removed at the coarse scales, since coarsening already
decorrelates and stride~3 would leave only 25 and 4 usable footprints.

Aggregating \texttt{randomError} deserves a note. The variance of a block mean
of $n$ footprints is $n^{-2}\sum_{ij}\rho_{ij}\sigma_i\sigma_j$, and under the
exponential correlation implied by \texttt{error\_corr\_cells} the $4\times4$
block retains $\sigma \times 0.744$ rather than the $\sigma/\sqrt{n} = \sigma
\times 0.25$ that independence would give. Assuming independence would have
overstated the satellite's precision by $3.0\times$ at $0.4^\circ$ and
$4.5\times$ at $0.8^\circ$.

The resulting arms carry comparable information: effective observation counts
(footprints divided by the correlation inflation) are $57$, $57$ and $50$ at
$0.1^\circ$, $0.4^\circ$ and $0.8^\circ$. The ladder therefore varies
observation \emph{scale} at nearly fixed observation \emph{quantity}, which was
not designed but makes the contrast much easier to interpret.

\subsection{Primary result: the satellite does not help}

Table~\ref{tab:ingestion-point} gives point skill at withheld gauges. The
quantity of interest is $\Delta = \mathrm{CRPS}(\text{combined}) -
\mathrm{CRPS}(\text{gauges})$, paired within the same dump file so that both
arms share the fold, the withheld stations, the days and the seeds. This is an
exact pairing rather than a fold-matched one, and it is the sharpest comparison
the design admits.

\begin{table}[t]
\centering
\caption{Point skill at 380 withheld station-days (178 wet). $\Delta$ is
combined minus gauges-only with a 95\% paired bootstrap interval; negative
favours adding the satellite. Background CRPS is $7.55$ for every arm, as it
must be.}
\label{tab:ingestion-point}
\small
\begin{tabular}{lrrrr@{\hspace{4pt}}l}
\hline
arm & gauges & satellite & combined & $\Delta$ & 95\% CI \\
\hline
S04   & 4.96 &  6.44 & \textbf{4.77} & $-0.192$ & $[-0.549, +0.112]$ \\
GW    & 4.95 &  6.48 & 4.82 & $-0.130$ & $[-0.433, +0.143]$ \\
RAW   & 4.96 &  6.48 & 4.85 & $-0.111$ & $[-0.413, +0.175]$ \\
SW    & 4.96 &  6.49 & 4.86 & $-0.102$ & $[-0.410, +0.185]$ \\
SL    & 4.96 &  6.89 & 4.93 & $-0.029$ & $[-0.289, +0.205]$ \\
RATIO & 4.95 &  6.89 & 4.94 & $-0.006$ & $[-0.271, +0.236]$ \\
S08   & 4.96 &  6.46 & 4.94 & $-0.016$ & $[-0.304, +0.246]$ \\
GL    & 4.99 &  6.48 & 4.95 & $-0.041$ & $[-0.357, +0.252]$ \\
MEASR & 4.98 &  6.56 & 5.01 & $+0.029$ & $[-0.261, +0.308]$ \\
GM    & 5.07 &  6.48 & 5.10 & $+0.030$ & $[-0.295, +0.344]$ \\
S1    & 4.96 & 13.23 & 7.35 & $+2.388$ & $[+1.824, +2.956]$ \\
\hline
\end{tabular}
\end{table}

\textbf{Ten of the eleven intervals contain zero.} Under the decision rule
fixed before the runs --- an arm counts as better only if its interval excludes
zero --- no configuration beats gauges-only, and gauges-only is the production
configuration on this evidence. The best point estimate, S04 at $-0.192$, is
not significant.

Every arm nonetheless improves substantially on the background ($7.55
\rightarrow 4.77$--$5.10$), so the assimilation is doing real work. The finding
is specifically that IMERG adds nothing detectable \emph{on top of} 30 gauges,
not that assimilation fails.

\subsection{Stride 1 at native resolution is actively harmful}

The single significant result is negative. S1 --- stride 1 at $0.1^\circ$,
identical to RAW in every other respect --- degrades CRPS by
$+2.388~[+1.824, +2.956]$, with satellite-arm CRPS of $13.23$ against
${\sim}6.5$ elsewhere, bias $+17.53$~mm~day$^{-1}$ and wet-area fraction
$0.990$.

The mechanism is the perturbation pathology above. At stride 1 on the
$0.1^\circ$ grid the automatic correlation inflation is $56.5\times$; because
perturbations are drawn with $\mathrm{sd} = \sqrt{R}$, they are ${\sim}7.5\times$
their nominal amplitude and the observation ceases to carry usable information.
This reproduces an earlier finding (\texttt{s1r10T}, $+2.60$~CRPS, May 2024) on
an independent window, so it should now be treated as established rather than
provisional.

\subsection{Secondary result: the satellite's usable information is at large scales}

Point skill is flat across the scale ladder, but spatial structure is not
(Table~\ref{tab:ingestion-structure}). Three metrics move monotonically with
footprint size, and they are not measuring the same thing.

\begin{table}[t]
\centering
\caption{Spatial structure by observation footprint. Pattern correlations are
daily, domain-wide, against each product. The background's best pattern
correlation is $0.60$.}
\label{tab:ingestion-structure}
\small
\begin{tabular}{llrrrrr}
\hline
arm & footprint & CHIRPS & CPC & IMERG & wet area & bias \\
\hline
background & ---              & ---    & ---   & (0.60) & 0.904 & $+8.37$ \\
S1         & $0.1^\circ$ str.\ 1 & $-0.078$ & 0.187 & 0.238 & 0.990 & $+17.53$ \\
RAW        & $0.1^\circ$ str.\ 3 & $-0.019$ & 0.415 & 0.499 & 0.922 & $+1.84$ \\
S04        & $0.4^\circ$      & $+0.011$ & 0.449 & 0.592 & 0.875 & $+2.16$ \\
S08        & $0.8^\circ$      & $+0.029$ & \textbf{0.538} & \textbf{0.673} & \textbf{0.820} & \textbf{$+1.19$} \\
\hline
\end{tabular}
\end{table}

The sharpest statement of the result is a comparison with the background rather
than between arms. \textbf{The background's own best pattern correlation is
$0.60$. Assimilating IMERG at $0.1^\circ$ reduces it to $0.50$; only the
$0.8^\circ$ arm improves on it, at $0.67$.} Assimilating the satellite at its
native footprint currently makes the spatial field worse than not assimilating
it at all. Coarsening removes retrieval noise the fine footprints were
injecting while retaining the large-scale constraint, which is precisely the
hypothesis the ladder was constructed to test.

Wet-area fraction and bias corroborate independently: $0.922 \rightarrow 0.875
\rightarrow 0.820$ and $+1.84 \rightarrow +2.16 \rightarrow +1.19$ across the
ladder, with the coarsest arm closest to plausibility on both.

\paragraph{Circularity, stated plainly.}
Pattern correlation against IMERG is partly self-fulfilling, since IMERG is
what is being assimilated. The CPC correlation is the cleaner signal --- CPC
does not enter the likelihood --- and it tells the same story ($0.415
\rightarrow 0.449 \rightarrow 0.538$), though CPC does condition the prior and
so is not fully independent either. Wet-area fraction and bias carry no such
circularity and agree. The result should be read as supported by three partly
independent lines rather than proven by any one.

\paragraph{The ladder is confounded with thinning, necessarily.}
RAW uses stride~3 while S04 and S08 use stride~1, because stride~3 at the
coarse scales would leave 25 and 4 footprints. The confound cannot be removed,
but its direction is known: S1 shows that stride~1 at $0.1^\circ$ is
catastrophic, so if thinning were driving the trend the coarse stride-1 arms
should be the \emph{worst}, and they are the best. The effective observation
counts (57/57/50) also match.

\subsection{What the strength axis showed}

Nothing separable, but with a consistent ordering worth recording. Combined
CRPS rises monotonically as the gauges are loosened: GW ($\sigma = 0.05$)
$4.82$, RAW ($0.10$) $4.85$, GL ($0.25$) $4.95$, GM ($0.41$) $5.10$. None of
these differences is significant.

MEASR --- which sets both streams to the representativeness measured from the
variogram, $\sigma_{\mathrm{rep}} = 0.410$ --- is among the worst arms
($5.01$, bias $+3.13$). This is worth stating carefully, because it is
counter-intuitive and easy to misread. The measured value is a correct
description of point-versus-cell mismatch; it is nevertheless not the best
operational choice \emph{given a background biased by $+8.37$~mm~day$^{-1}$},
because loosening the gauges admits more of that bias into the analysis. When
the prior is badly biased, the statistically correct observation error and the
skill-optimal observation error are different quantities. The correct
conclusion is not that the measurement was wrong but that the prior is.

\subsection{What remains broken}

\begin{enumerate}
  \item \textbf{Every arm is outside the product wet-area envelope}, best case
    $0.820$. The background is $0.904$ with a bias of $+8.37$~mm~day$^{-1}$ at
    withheld stations. Assimilation removes most of the bias
    ($+8.37 \rightarrow +1.19$ at best) but cannot fix a prior that rains
    almost everywhere.
  \item \textbf{The analysis has essentially no daily agreement with CHIRPS}
    ($-0.078$ to $+0.029$ across arms) --- the product the prior was
    \emph{trained on}. Combined with CHIRPS's measured daily correlation of
    $0.29$--$0.56$ against these same gauges, against IMERG's $0.71$--$0.78$,
    this points at the training target rather than at any assimilation setting.
  \item \textbf{IMERG has still never been bias-corrected.}
    \texttt{observations.imerg.bias\_correction} exists in the configuration
    and is not read by the assimilation script, so every result above uses raw
    V07B with a measured $+5.56$~mm~day$^{-1}$ bias at the gauges, in a
    Gaussian likelihood that assumes unbiasedness.
\end{enumerate}

\subsection{Power, and what would settle the scale question}

The point-skill test is underpowered for the effect sizes observed. S04's
interval half-width is $0.33$~mm~day$^{-1}$ against a point estimate of
$0.192$; excluding zero requires roughly $(0.33/0.192)^2 \approx 3.0$ times the
sample, that is ${\sim}1130$ withheld station-days, or about 30 days rather
than 10.

The confirmation run is therefore the full month of May 2022 (31 days, five
folds, ${\sim}1240$ withheld station-days) restricted to RAW, S04 and S08 ---
the scale ladder only, since the strength axis produced nothing to follow up
and S1 is established as harmful.

One caveat on that design: the 31-day window is a strict \emph{superset} of the
10-day window, so it increases power rather than providing an independent
replication. If the scale result survives, June 2022 gives a genuinely
independent test on data not used here.
